\begin{example}[Refuted: a positive uniform lower bound for fierce excess]
\label{D:fierce-small-excess}
The assertion that Galois extensions with inseparable residue have a
positive uniform lower bound for $(\delta-e+1)/e$ is false.
We fix an algebraically closed field $k$ of characteristic $p>0$.
For every power $q=p^r>1$, there is a finite Galois extension of
two-variable rational function fields over $k$, with a divisorial valuation
having ramification index $e=q(q-1)$, purely inseparable residue degree $q$,
and different exponent
\[
  \delta=q^2+q-3.
\]
Consequently
\[
  \frac{\delta-e+1}{e}=\frac2q\longrightarrow0.
\]
In particular, inseparability of the residue extension does not give a
positive lower bound for this normalized excess different, even at a fixed
characteristic and with a Galois extension between function fields of smooth
surfaces.  This is a local ramification example; we do not assert that it
extends to a quasi-\'etale cover of a prescribed contracted surface.
\end{example}

\begin{proof}
We take algebraically independent elements $s,z$ over $k$, put
$t=s^{q-1}$, and adjoin elements $w,y$ satisfying
\[
 w^q-w=s^{-1},\qquad y^q-y=zs^{-1}.
\]
We write $K=k(t,z)$, $T=k(s,z)$, and $L=T(w,y)$.
The extension $T/K$ has degree $q-1$ and is totally tamely ramified
at $t=0$.  For $\pi=w^{-1}$ the first equation becomes
\[
 \pi^q+s\pi^{q-1}-s=0.
\]
This polynomial is Eisenstein at $s=0$.  Thus $[T(w):T]=q$, the
ramification index is $q$, and its residue field is $k(z)$.
Putting $u=y/w$, we obtain
\begin{equation}
\label{D:fierce-small-excess-coordinates}
 s=\frac{\pi^q}{1-\pi^{q-1}},\qquad
 z=\frac{u^q-\pi^{q-1}u}{1-\pi^{q-1}},\qquad
 L=k(\pi,u).
\end{equation}
The monic equation
\[
 u^q-\pi^{q-1}u-z(1-\pi^{q-1})=0
\]
reduces to $u^q-z$ at $\pi=0$.  This is irreducible over $k(z)$,
since the transcendental element $z$ is not a $p$th power.
The corresponding finite algebra over the discrete valuation ring
$k(z)[\pi]_{(\pi)}$ is local with maximal ideal generated by $\pi$;
it is therefore a discrete valuation ring.  This proves
$[L:T(w)]=q$, ramification index one, and purely inseparable residue
degree $q$.  We have proved $[L:K]=q^2(q-1)$.

For $a,b\in\mathbf F_q$ and $c\in\mathbf F_q^\times$ the formulas
\[
 (w,y,s,z)\longmapsto(cw+a,cy+b,c^{-1}s,z)
\]
give $q^2(q-1)$ distinct automorphisms fixing $t$ and $z$.
They form $\mathbf F_q^2\rtimes\mathbf F_q^\times$, with the
scalar action on the two additive factors.  The extension $L/K$ is
therefore Galois.  Each automorphism preserves the valuation at $\pi=0$:
its action on the displayed coordinates is
\[
 \pi\longmapsto\frac{\pi}{c+a\pi},\qquad
 u\longmapsto\frac{cu+b\pi}{c+a\pi}.
\]
In particular it acts trivially on the residue field $k(u)$.
The ramification index and residue degree over $K$ are consequently
$e=q(q-1)$ and $f=q$, respectively, as also follows by multiplying
the indices in the three steps above.

The different at this divisorial valuation is the order of the
Jacobian determinant for $(\pi,u)\mapsto(t,z)$.
Writing $D=1-\pi^{q-1}$, differentiation in characteristic $p$ gives
\[
 \frac{\partial s}{\partial\pi}
       =-\frac{\pi^{2q-2}}{D^2},\qquad
 \frac{\partial t}{\partial\pi}
       =\frac{\pi^{q^2-2}}{D^q},\qquad
 \frac{\partial z}{\partial u}
       =-\frac{\pi^{q-1}}D.
\]
Since $t$ is independent of $u$, we conclude
\[
 dt\wedge dz
   =-\frac{\pi^{q^2+q-3}}{D^{q+1}}\,d\pi\wedge du.
\]
The factor $D$ is a unit at the valuation, so
$\delta=q^2+q-3$ and the stated normalized excess follows.
\end{proof}

\begin{remark}[What this example leaves open]
In the Galois case, Proposition~\ref{D:single-curve-residue-tameness}
gives the necessary bound
\[
 \frac{\delta-e+1}{e}
   \leq\frac{2g-2}{d}\left(1-\frac1{f_{\mathrm{insep}}}\right).
\]
For every fixed $g>1$ and $d>0$, the numerical ramification data of
Example~\ref{D:fierce-small-excess} satisfy this inequality for all
sufficiently large $q$.  Thus the inequality by itself cannot exclude
fierce residue extensions by imposing a large value of $d$.
Realizability over a contracted surface, and the additional global
conditions in the finite-generation question, remain separate issues.
\end{remark}

\paragraph{Finite certificate.}
The program \nolinkurl{parts/C/code/check_fierce_different.py} and its
output \nolinkurl{parts/C/output/fierce_different_certificate.json}
check $q=2,3,4,5,8,9$.  They verify the finite fields, the faithful affine
semidirect actions and their generation, the two defining equations, all
three derivatives, the full Jacobian identity, and its valuation.
The run passed $1{,}029{,}677$ exact assertions in
$0.551970493994$ seconds.  The proof for every prime power is the argument
above, not an extrapolation from these finite checks.
